% !TeX root = ../../Book3.tex
% 纲要标题：## 第4章　关联素理想与准素分解
% 纲要位置：工作记录/计划纲要.md，第 2005 行。

\chapter{关联素理想与准素分解}
\label{b3:ch:04}
\BookChapterNumber{b3:ch:04}

\begin{chapterabstract}
本章用模中元素的零化子定义关联素理想，证明其存在性、有限性与局部化公式，并由此描述零因子。随后证明 Noether 环中的理想，以及这类环上有限生成模中的子模，均存在准素分解，并区分孤立准素分量与嵌入准素分量的唯一性。
\end{chapterabstract}

\begin{learninggoals}
\begin{itemize}
\item 由实际元素的零化子寻找关联素理想，说明存在性与有限性各需什么假设。
\item 用素理想滤过和局部化研究有限生成模，以关联素理想的并描述零因子，并辨认支集中的极小点。
\item 证明准素分解存在性，准确表述两类唯一性，给出素理想幂不准素及嵌入准素分量不唯一的例子。
\item 对理想、子模与矩阵余核作准素分析，区分非约化结构与嵌入关联素理想。
\end{itemize}
\end{learninggoals}

设 \(k\) 为域。在 \(R=k[x,y]/(xy)\) 中，\(\bar x\bar y=0\)，因而两个非零元素都是零因子。这里可以直接从两条分支理解：\(\bar x\) 在 \(y\) 轴上消失，\(\bar y\) 在 \(x\) 轴上消失。但取 \(I=(x^2,xy)\)、\(A=k[x,y]/I\)，有 \(\sqrt I=(x)\)，所以 \(A\) 作为 \(k[x,y]\) 模的支集只是 \(V(x)\)。同时，\(A\) 中的非零元 \(\bar x\) 满足 \(x\bar x=y\bar x=0\)，其零化子是 \((x,y)\)。约化商 \(A/(\bar x)\cong k[y]\) 有同一个底闭集，却没有这份原点处的额外零化信息；支集本身不能辨认这种区别。

\ProseParagraphBreak

要找出支集没有区分的信息，应当反过来问：哪些标量杀死了给定的非零模元素？得到的零化子中若有素理想，便形成关联素理想。它们不仅标出不可约分支的一般点，也可能检测到分支内部特殊点处的额外模结构。准素分解把这些素理想及其相应的商模结构落实为理想的交表示；随后还须分清哪些素理想由原理想确定，哪些分量会随分解的选择改变。

\ProseParagraphBreak

本章使用\chapref{b3:ch:01}的素谱、根理想及有限素理想回避，\chapref{b3:ch:02}的模局部化和支集，以及\chapref{b3:ch:03}的 Noether 模与有限性。理想商的基础运算可结合第一卷第10.3节复习。

\ProseParagraphBreak
除这条带有嵌入信息的直线外，环 \(k[x,y,z]/(xy-z^2)\) 还提供了另一项检验：即使环境是整环，素理想的普通幂也未必准素。它要求我们直接检查准素定义中的乘法条件，不能只从理想的根推测准素性。

% !TeX root = ../../Book3.tex
% 纲要标题：### 4.1 关联素理想
% 纲要位置：工作记录/计划纲要.md，第 2007 行。
% 纲要细目（写作范围）：
% * 元素零化子的素性；循环子模与素商模的嵌入、短正合列的非对称性
% * 关联素理想的存在和有限性；非典范素理想滤过、局部化中的统一清分母
% * 零因子作为关联素理想的并；支集极小点与孤立、嵌入关联素理想

\section{关联素理想}
\label{b3:sec:0401}

同一个元素可以被许多环元素零化，但它的全部零化子组成一个确定的理想。对交换含幺环 \(R\) 的模 \(M\) 及元素 \(m\in M\)，映射 \(a\mapsto am\) 给出 \(R\) 模同构 \(Rm\cong R/\operatorname{Ann}_R(m)\)。当这个零化子是素理想 \(\mathfrak p\) 时，模中便含有一个与 \(R/\mathfrak p\) 同构的循环子模。这给出了从模中提取素理想的方法。在 Noether 环上，每个非零模都有这样的素理想，有限生成模的关联素理想还只有有限个。前一个结论不要求模有限生成，而这些假设也不是每个具体模存在关联素理想的必要条件。

\subsection{元素零化子的素性}
\label{b3:subsec:0401:01}

\begin{definition}\label{b3:def:associated-prime}
设 \(R\) 为交换含幺环，\(M\) 为 \(R\) 模，\(m\in M\)。元素 \(m\) 的零化子为
\[
\operatorname{Ann}_R(m)=\{\,a\in R\mid am=0\,\}.
\]
若 \(R\) 的素理想 \(\mathfrak p\) 等于 \(M\) 中某个非零元素的零化子，则称 \(\mathfrak p\) 为 \(M\) 的\term[guanlian-sulixiang]{关联素理想}[associated prime]。所有这样的素理想组成集合 \(\operatorname{Ass}_R(M)\)，并约定 \(\operatorname{Ass}_R(0)=\varnothing\)。
\symidx{\(\operatorname{Ass}_R(M)\)：模的关联素理想集合}
\end{definition}

因此 \(\mathfrak p\in\operatorname{Ass}_R(M)\) 等价于存在 \(R\) 模单射 \(R/\mathfrak p\hookrightarrow M\)。定义说的是某个元素的零化子，不是整个模的零化子必须为素理想。

\begin{lemma}\label{b3:lem:maximal-element-annihilator-prime}
设 \(R\) 为交换含幺环，\(M\) 为 \(R\) 模，\(0\ne m\in M\)。若 \(\operatorname{Ann}_R(m)\) 在集合 \(\{\operatorname{Ann}_R(x):0\ne x\in M\}\) 中按包含关系极大，则它是 \(R\) 的素理想。
\end{lemma}
\begin{proof}
记 \(I=\operatorname{Ann}_R(m)\)。因 \(m\ne0\)，\(I\ne R\)。若 \(ab\in I\) 且 \(b\notin I\)，则 \(bm\ne0\)。乘以 \(b\) 后，原来零化 \(m\) 的元素仍零化 \(bm\)，所以 \(I\subseteq\operatorname{Ann}_R(bm)\)。由于 \(bm\) 仍是非零元素，极大性不允许这个零化子严格增大，只能有 \(\operatorname{Ann}_R(bm)=I\)。此时 \(a(bm)=0\) 说明 \(a\) 原来就属于 \(I\)，这正是素性。
\end{proof}

这里的“极大”指元素零化子这一族中的极大元，并不表示 \(I\) 是环的极大理想。例如整环 \(R\) 作为自身上的模时，每个非零元素的零化子都是 \((0)\)，于是 \(\operatorname{Ass}_R(R)=\bigl\{(0)\bigr\}\)，即使 \(R\) 不是域也成立。

\begin{proposition}\label{b3:prop:associated-exact-sequence}
设 \(R\) 为交换含幺环，\(0\longrightarrow N\longrightarrow M\longrightarrow P\longrightarrow0\) 为 \(R\) 模的短正合列。则
\[
\operatorname{Ass}_R(N)\subseteq\operatorname{Ass}_R(M)
 \subseteq\operatorname{Ass}_R(N)\cup\operatorname{Ass}_R(P).
\]
特别地，有限直和的关联素理想集合是各因子对应集合的并。对一般短正合列，却不能保证 \(\operatorname{Ass}_R(P)\subseteq\operatorname{Ass}_R(M)\)。
\end{proposition}
\begin{proof}
子模中的元素在 \(N\) 与 \(M\) 内有同一个零化子，得第一个包含。给定 \(\mathfrak p\in\operatorname{Ass}_R(M)\)，在 \(M\) 内取 \(L\cong R/\mathfrak p\)。\(L\) 中每个非零元素的零化子都是 \(\mathfrak p\)，因为 \(R/\mathfrak p\) 是整环。如果 \(L\cap N\ne0\)，取交内非零元素即得 \(\mathfrak p\in\operatorname{Ass}_R(N)\)；如果交为零，\(L\) 单射到 \(P\)，故 \(\mathfrak p\in\operatorname{Ass}_R(P)\)。直和结论由分裂正合列和两因子各自的嵌入得到，再对因子数归纳。

商模可以产生原模没有的关联素理想。例如，设 \(k\) 为域，\(R=k[t]\)，考虑短正合列
\[
0\longrightarrow tR\longrightarrow R\longrightarrow R/(t)\longrightarrow0.
\]
环 \(R\) 是整环，故 \(\operatorname{Ass}_R(R)=\{(0)\}\)；而 \(R/(t)\cong k\) 的每个非零元素的零化子均为 \((t)\)，所以 \(\operatorname{Ass}_R(R/(t))=\{(t)\}\)。因此商模的关联素理想不必属于中间模的关联素理想集合。
\end{proof}

\subsection{关联素理想的存在性与有限性}
\label{b3:subsec:0401:02}

\begin{theorem}\label{b3:thm:associated-existence}
若 \(R\) 是交换含幺 Noether 环，\(M\ne0\) 是 \(R\) 模，则 \(\operatorname{Ass}_R(M)\ne\varnothing\)。这里不要求 \(M\) 有限生成。
\end{theorem}
\begin{proof}
所有非零元素零化子构成非空理想族。Noether 升链条件保证其中有极大元，再用\cref{b3:lem:maximal-element-annihilator-prime}。
\end{proof}

Noether 环上的有限生成模可以逐层拆成素商模，再由短正合列的包含关系控制它的关联素理想。每一步的素商子模需要作出选择，所得滤过通常不唯一，其中出现的素理想只给关联素理想集合一个上界。

\begin{theorem}\label{b3:thm:associated-finite}
设 \(R\) 是交换含幺 Noether 环，\(M\) 是有限生成 \(R\) 模。存在有限子模链
\[
0=M_0\subsetneq M_1\subsetneq\cdots\subsetneq M_r=M,
\qquad M_i/M_{i-1}\cong R/\mathfrak p_i,
\]
其中各 \(\mathfrak p_i\) 为素理想；这种链称为\term[sulixiang-lvguo]{素理想滤过}[prime filtration]。并且
\[
\operatorname{Ass}_R(M)\subseteq\{\mathfrak p_1,\ldots,\mathfrak p_r\},
\]
所以关联素理想只有有限个。\(M=0\) 时采用空滤过。
\end{theorem}
\begin{proof}
若已构造 \(M_i\ne M\)，对非零商模 \(M/M_i\) 应用\cref{b3:thm:associated-existence}，取其中同构于 \(R/\mathfrak p_{i+1}\) 的子模，并令 \(M_{i+1}\) 为其原像。由\cref{b3:prop:noether-finite-modules}，有限生成模 \(M\) 是 Noether 模，故这个严格上升过程必在有限步停止。对短正合列
\(0\rightarrow M_{i-1}\rightarrow M_i\rightarrow R/\mathfrak p_i\rightarrow0\)
逐次应用\cref{b3:prop:associated-exact-sequence}，再用 \(\operatorname{Ass}_R(R/\mathfrak p_i)=\{\mathfrak p_i\}\)，得到所需包含。
\end{proof}

例如 \(\mathbb Z\supset2\mathbb Z\supset0\) 的两个因子为 \(\mathbb Z/2\mathbb Z\) 与 \(\mathbb Z\)，但 \(\operatorname{Ass}_{\mathbb Z}(\mathbb Z)=\bigl\{(0)\bigr\}\)，滤过中的 \((2)\) 是额外出现的素理想。此外，有限生成性不能删除：\(\bigoplus_{p\;\text{为素数}}\mathbb Z/p\mathbb Z\) 的关联素理想包括所有 \((p)\)，因而是无限集。

\begin{proposition}\label{b3:prop:associated-localization}
设 \(R\) 为交换含幺 Noether 环，\(S\subseteq R\) 为含一的乘法集，\(M\) 为任意 \(R\) 模。则
\[
\operatorname{Ass}_{S^{-1}R}(S^{-1}M)
=\bigl\{\,S^{-1}\mathfrak p\mid \mathfrak p\in\operatorname{Ass}_R(M),\quad
                       \mathfrak p\cap S=\varnothing\,\bigr\}.
\]
\end{proposition}
\begin{proof}
若 \(\operatorname{Ann}_R(m)=\mathfrak p\) 且 \(\mathfrak p\cap S=\varnothing\)，则 \(m/1\ne0\)。当 \((a/s)(m/1)=0\) 时，有 \(t\in S\) 使 \(tam=0\)；由 \(ta\in\mathfrak p\)、\(t\notin\mathfrak p\) 得 \(a\in\mathfrak p\)。故其零化子恰为 \(S^{-1}\mathfrak p\)。

反之，设 \(\mathfrak q\) 是非零元素 \(m/s\) 的素零化子，令 \(\mathfrak p\) 为 \(\mathfrak q\) 在 \(R\) 中的收缩。由\cref{b3:thm:spec-localization}，\(\mathfrak p\) 是避开 \(S\) 的素理想，且 \(\mathfrak q=S^{-1}\mathfrak p\)。以单位 \(s/1\) 乘此元素后可改用 \(m/1\)。要在原模中找到零化子恰为 \(\mathfrak p\) 的元素，必须把局部成立的零化关系同时变成原模中的等式；这一步需要共同的分母。由于 \(R\) 是 Noether 环，理想 \(\mathfrak p\) 有有限生成集，写成 \(\mathfrak p=(a_1,\ldots,a_l)\)。因 \((a_i m)/1=0\)，可选 \(s_i\in S\) 使 \(s_i a_i m=0\)。令 \(t=s_1\cdots s_l\)，此时 \(\mathfrak p\subseteq\operatorname{Ann}_R(tm)\)；当 \(\mathfrak p=(0)\) 时采用空生成集并令 \(t=1\)。若 \(a(tm)=0\)，由于 \(t/1\) 可逆，局部化后 \(a/1\) 便零化 \(m/1\)，所以 \(a\in\mathfrak p\)。于是 \(\operatorname{Ann}_R(tm)=\mathfrak p\)，且 \(tm\ne0\)，否则 \(m/1=0\)。这就找到了原模中的关联元素，而所用的有限生成性来自理想 \(\mathfrak p\)，并未要求 \(M\) 有限生成。
\end{proof}

例如，设 \(k\) 为域，\(R=k[x,y]\)、\(M=R/(x)\)、\(\mathfrak m=(x,y)\)。因 \(R/(x)\cong k[y]\) 是整环，每个非零模元素的零化子都是 \((x)\)，故 \(\operatorname{Ass}_R(M)=\{(x)\}\)。局部化后有
\[
\begin{aligned}
M_{\mathfrak m}&\cong R_{\mathfrak m}/(x)R_{\mathfrak m}
\cong k[y]_{(y)},\\
\operatorname{Ass}_{R_{\mathfrak m}}(M_{\mathfrak m})
&=\{(x)R_{\mathfrak m}\}.
\end{aligned}
\]
这个唯一的关联素理想严格小于局部环的极大理想 \(\mathfrak mR_{\mathfrak m}\)：元素 \(y/1\) 属于后者，在商环 \(k[y]_{(y)}\) 中的像却非零，所以不属于前者。

\begin{proposition}\label{b3:prop:minimal-primes-associated}
设 \(R\) 为交换含幺 Noether 环，\(M\) 为非零有限生成 \(R\) 模。则 \(\operatorname{Supp}_R(M)\) 的每个素理想都包含某个关联素理想，且
\[
\min\operatorname{Supp}_R(M)=\min\operatorname{Ass}_R(M).
\]
尤其，对任意真理想 \(I\)，包含 \(I\) 的极小素理想都是 \(R/I\) 的关联素理想。
\end{proposition}
\begin{proof}
若 \(\mathfrak p\in\operatorname{Ass}_R(M)\)，上述局部化命题保证 \(M_{\mathfrak p}\ne0\)，故 \(\operatorname{Ass}_R(M)\subseteq\operatorname{Supp}_R(M)\)。给定 \(\mathfrak q\in\operatorname{Supp}_R(M)\)，\(R_{\mathfrak q}\) 是 Noether 环，非零模 \(M_{\mathfrak q}\) 有关联素理想；由\cref{b3:prop:associated-localization}，它来自某个 \(\mathfrak p\in\operatorname{Ass}_R(M)\)，且 \(\mathfrak p\subseteq\mathfrak q\)。因此两集合的极小元一致。对有限生成模，\cref{b3:thm:support-finite}给出 \(\operatorname{Supp}_R(M)=V\bigl(\operatorname{Ann}_R(M)\bigr)\)。取 \(M=R/I\) 就得到最后一句。
\end{proof}

设 \(k\) 为域，取 \(R=k[x,y]\)、\(M=R/(x^2,xy)\)。由 \(\sqrt{(x^2,xy)}=(x)\)，其支集为 \(V(x)\)，极小元只有 \((x)\)，所以 \((x)\) 也是极小关联素理想。另一方面，非零元素 \(\bar x\) 的零化子为 \((x,y)\)，具体计算见\subsecref{b3:subsec:0404:02}。因此 \((x,y)\) 也是关联素理想，却严格包含 \((x)\)。

\subsection{零因子的关联素理想描述}
\label{b3:subsec:0401:03}

设 \(R\) 为交换含幺环，\(M\) 为 \(R\) 模。若存在 \(0\ne m\in M\) 使 \(am=0\)，则称 \(a\in R\) 是 \(M\) 上的\term[mo-shang-lingyinzi]{零因子}[zero divisor on a module]；否则称它在 \(M\) 上是正则的。记这些零因子的集合为 \(Z_R(M)\)。这里允许 \(a=0\)：对非零模，零当然属于这个集合；对零模，集合为空。

\ProseParagraphBreak
标量 \(a\) 的乘法映射 \(M\to M\)、\(m\mapsto am\) 单射，当且仅当 \(a\notin Z_R(M)\)。因此零因子集合决定哪些标量在模上的作用失去单射性；在 Noether 环上，关联素理想恰好给出这个集合的描述。

\begin{theorem}\label{b3:thm:associated-zero-divisors}
设 \(R\) 为交换含幺 Noether 环，\(M\) 为 \(R\) 模。则
\[
Z_R(M)=\bigcup_{\mathfrak p\in\operatorname{Ass}_R(M)}\mathfrak p.
\]
若 \(M\) 有限生成，这就是有限个素理想的并。
\end{theorem}
\begin{proof}
关联素理想中的元素会零化定义它的非零元素，因此右端包含于左端。反向设 \(am=0\)、\(m\ne0\)。包含 \(\operatorname{Ann}_R(m)\) 的非零元素零化子构成非空理想族；由于 \(R\) 是 Noether 环，其中存在极大元 \(\operatorname{Ann}_R(n)\)。它在全部非零元素零化子中也是极大的，因任何更大零化子仍包含 \(\operatorname{Ann}_R(m)\)。\cref{b3:lem:maximal-element-annihilator-prime}说明它是关联素理想，而 \(a\) 在其中。
\end{proof}

例如，设 \(k\) 为域，在 \(R=k[x,y]/(xy)\) 中，\(\bar x\) 与 \(\bar y\) 的零化子分别为 \((\bar y)\) 和 \((\bar x)\)。每个元素都可唯一写为常数加一个无常数项的 \(x\) 多项式，再加一个无常数项的 \(y\) 多项式。限制到两轴给出单射 \(R\to k[x]\times k[y]\)；不在 \((\bar x)\cup(\bar y)\) 中的元素在两个坐标上都非零，故乘积为零时，两个整环中的消去律迫使另一乘子为零。而 \((\bar x)\)、\((\bar y)\) 中的元素分别被 \(\bar y\)、\(\bar x\) 零化，所以零因子恰为这两个素理想的并。\(\bar x+\bar y\) 不在其中，因而不是零因子，尽管它的两个加项都是零因子。这说明零因子集合不一定对加法封闭，因而不一定是理想。

\ProseParagraphBreak
更一般地，设 \(R\) 为交换含幺 Noether 环，若 \(R\) 的理想 \(J\) 的每个元素都是有限生成 \(R\) 模 \(M\) 上的零因子，由\cref{b3:thm:associated-finite}和\cref{b3:thm:associated-zero-divisors}及\cref{b3:lem:finite-prime-avoidance}的有限素理想回避性质，\(J\) 必包含于某一个关联素理想。这把“逐个检查乘法是否单射”的问题转成了有限个素理想的包含问题。以上存在性、局部化和滤过方法可对照 \cite[\S6, Theorems 6.1--6.5]{Matsumura1989CommutativeRingTheory}。

\ProseParagraphBreak
设 \(R\) 为交换含幺 Noether 环，\(M\) 为有限生成 \(R\) 模。按包含关系极小的关联素理想称为\term[guli-guanlian-sulixiang]{孤立关联素理想}[isolated associated prime]，其余称为\term[qianru-guanlian-sulixiang]{嵌入关联素理想}[embedded associated prime]。这里的“孤立”指素理想包含关系中的极小性，不要求它是 Zariski 拓扑中的孤立点，即不要求单点集合为开集。由\cref{b3:prop:minimal-primes-associated}，前者就是支集的极小元，后者则严格包含某个极小关联素理想；两类都由 \(M\) 本身确定。\cref{b3:tab:support-associated-primes}汇集了这些关系。

\begin{table}[htbp]
\centering
\caption[支集与关联素理想]{支集与关联素理想。设 \(R\) 为交换含幺 Noether 环，\(M\) 为有限生成 \(R\) 模；极小元均按素理想的包含关系理解。}
\label{b3:tab:support-associated-primes}
\begin{tabularx}{\textwidth}{@{}>{\raggedright\arraybackslash}p{0.25\textwidth}>{\raggedright\arraybackslash}X@{}}
\toprule
集合 & 与支集及零化子的关系 \\
\midrule
支集 & \(\operatorname{Supp}_R(M)=V\bigl(\operatorname{Ann}_R(M)\bigr)\)；每个支集中的素理想都包含某个关联素理想 \\
关联素理想 & \(\operatorname{Ass}_R(M)\subseteq\operatorname{Supp}_R(M)\)，且该集合有限 \\
孤立关联素理想 & \(\min\operatorname{Ass}_R(M)=\min\operatorname{Supp}_R(M)\)，即支集的极小点 \\
嵌入关联素理想 & \(\operatorname{Ass}_R(M)\setminus\min\operatorname{Supp}_R(M)\)，即支集中非极小的关联点 \\
\bottomrule
\end{tabularx}
\end{table}

% !TeX root = ../../Book3.tex
% 纲要标题：### 4.2 准素理想
% 纲要位置：工作记录/计划纲要.md，第 1913 行。
% 纲要细目（写作范围）：
% * 准素性与根
% * 素幂不必准素的注意事项
% * 局部化中的准素理想

\section{准素理想}
\label{b3:sec:0402}

素理想在商环中排除一切非零零因子，准素理想则允许零因子存在，但要求它们最终被某个幂消去。这个差别使准素理想既保留一个确定的素支集，又能描述支集上的幂零结构。先把准素性直接写成乘法条件，再研究它与局部化的相容性。

\subsection{准素性与根理想}
\label{b3:subsec:0402:01}

\begin{definition}\label{b3:def:primary-ideal}
设 \(R\) 为交换含幺环，\(Q\subsetneq R\) 为真理想。若对任意 \(a,b\in R\) 都有
\[
ab\in Q,\quad b\notin Q
\quad\Longrightarrow\quad a^n\in Q\ \text{对某个整数 }n\ge1.
\]
则称 \(Q\) 为\term[zhunsu-lixiang]{准素理想}[primary ideal]。若 \(\sqrt Q=\mathfrak p\)，也称 \(Q\) 是 \(\mathfrak p\) 准素理想。
\end{definition}

等价地，\(R/Q\) 中每个零因子都是幂零元。素理想满足更强的条件 \(n=1\)，因而总是准素理想；反之不成立，例如 \(\mathbb Z/(p^e)\) 中每个非单位都是 \(p\) 的倍数，故 \((p^e)\) 为 \((p)\) 准素理想。

\begin{proposition}\label{b3:prop:primary-radical-colon}
设 \(R\) 为交换含幺环，\(Q\) 为 \(R\) 的准素理想。则其根 \(\mathfrak p=\sqrt Q\) 是素理想。对 \(c\notin Q\)，理想商 \((Q:c)\) 仍为 \(\mathfrak p\) 准素理想；若还满足 \(c\notin\mathfrak p\)，则 \((Q:c)=Q\)。同一素理想所属的有限个准素理想的交仍为该素理想所属的准素理想。
\end{proposition}
\begin{proof}
若 \(ab\in\sqrt Q\)、\(b\notin\sqrt Q\)，则某个 \(a^m b^m\in Q\)，且 \(b^m\notin Q\)。准素性给出 \(a^{mn}\in Q\)，从而 \(a\in\sqrt Q\)。根是真理想，故为素理想。

若 \(ab\in(Q:c)\)、\(b\notin(Q:c)\)，即 \(a(bc)\in Q\)、\(bc\notin Q\)，则某个 \(a^n\in Q\subseteq(Q:c)\)。又 \(Q\subseteq(Q:c)\)，而 \(dc\in Q\)、\(c\notin Q\) 蕴涵 \(d\in\mathfrak p\)，所以两者根相同。若 \(c\notin\mathfrak p\)，从 \(dc\in Q\) 与 \(d\notin Q\) 将推出 \(c\in\mathfrak p\)，故 \(d\in Q\)。

最后令 \(Q=\bigcap_{i=1}^rQ_i\)，各 \(Q_i\) 的根均为 \(\mathfrak p\)。若 \(ab\in Q\)、\(b\notin Q\)，则至少一个因子给出 \(a\in\mathfrak p\)。于是每个 \(Q_i\) 都包含 \(a\) 的某个幂，取这些指数的最大值即得 \(a^n\in Q\)；同时 \(\sqrt Q=\bigcap_i\sqrt{Q_i}=\mathfrak p\)。
\end{proof}

\begin{proposition}\label{b3:prop:maximal-radical-primary}
设 \(R\) 为交换含幺环，\(Q\) 为 \(R\) 的真理想。若其根是极大理想 \(\mathfrak m\)，则 \(Q\) 是 \(\mathfrak m\) 准素理想。因此极大理想的任意正整数次幂均为准素理想。
\end{proposition}
\begin{proof}
\(R/Q\) 只有一个极大理想 \(\mathfrak m/Q\)。若 \(a\notin\mathfrak m\)，其像为单位，故 \(ab\in Q\) 必导致 \(b\in Q\)。取逆否命题即得到准素条件。又 \(\sqrt{\mathfrak m^n}=\mathfrak m\)，得最后一句。
\end{proof}

\begin{proposition}\label{b3:prop:primary-associated-singleton}
设 \(R\) 为交换含幺 Noether 环，\(Q\subsetneq R\) 为真理想。则 \(Q\) 是准素理想，当且仅当 \(\operatorname{Ass}_R(R/Q)\) 仅有一个元素；这唯一的元素就是 \(\sqrt Q\)。
\end{proposition}
\begin{proof}
若 \(Q\) 为准素理想，对每个非零 \(\bar b\in R/Q\)，\(\operatorname{Ann}_R(\bar b)=(Q:b)\) 的根由\cref{b3:prop:primary-radical-colon}等于 \(\sqrt Q\)。因此当零化子本身为素理想时，它只能是 \(\sqrt Q\)；存在性保证这样的零化子确实出现。

反之，设关联素理想只有 \(\mathfrak p\)。\cref{b3:prop:minimal-primes-associated}说明 \(Q\) 的极小素理想只有 \(\mathfrak p\)，故 \(\sqrt Q=\mathfrak p\)。由\cref{b3:thm:associated-zero-divisors}，商环中的每个零因子来自 \(\mathfrak p\)，所以它都是幂零元，\(Q\) 准素。
\end{proof}

\subsection{素理想幂不为准素理想的例子}
\label{b3:subsec:0402:02}

极大理想幂的结论不能直接推广到任意素理想。以下例子甚至发生在一个有限型整环中。

\begin{example}\label{b3:exm:prime-square-not-primary}
设 \(k\) 为任意域，
\[
R=k[x,y,z]/(xy-z^2),\qquad \mathfrak p=(\bar x,\bar z).
\]
则 \(R\) 是整环，\(\mathfrak p\) 是素理想，而 \(\mathfrak p^2\) 不是准素理想。
\end{example}
\begin{proof}
映射 \(\bar x\mapsto s^2\)、\(\bar y\mapsto t^2\)、\(\bar z\mapsto st\) 将 \(R\) 嵌入 \(k[s,t]\)。具体地，利用 \(z^2=xy\)，每个类可写为 \(a(x,y)+z\cdot b(x,y)\)；代入后的第一项只含偶数指数对，第二项只含奇数指数对，两类单项式不能抵消，因此映射单射，\(R\) 是整环。又 \(R/\mathfrak p\cong k[y]\)，故 \(\mathfrak p\) 素。

关系 \(xy-z^2\) 是全次数二的齐次式，因而商环仍按全次数分解，互不相同的次数不能相互抵消。理想 \(\mathfrak p^2=(\bar x^2,\bar x\bar z,\bar z^2)\) 由全次数二的齐次元生成，所以非零的一次齐次元 \(\bar x\) 不属于它。然而
\[
\bar y\bar x=\bar z^2\in\mathfrak p^2,
\qquad \bar y^n\notin\mathfrak p\supseteq\mathfrak p^2\quad(n\ge1),
\]
后一句由商环 \(k[y]\) 核验。这直接违反准素条件。
\end{proof}

失效的原因是 \(\bar y\) 虽不在素理想 \(\mathfrak p\) 中，却在 \(R/\mathfrak p^2\) 中成为零因子。只知道 \(\sqrt{\mathfrak p^2}=\mathfrak p\) 并不足以推出准素性；上一小节中的“根是极大理想”确有作用。

\subsection{局部化中的准素理想}
\label{b3:subsec:0402:03}

\begin{proposition}\label{b3:prop:primary-localization}
设 \(R\) 为交换含幺环，\(Q\) 为 \(R\) 的 \(\mathfrak p\) 准素理想，\(S\subseteq R\) 为含一的乘法集。若 \(S\cap\mathfrak p\ne\varnothing\)，则 \(S^{-1}Q=S^{-1}R\)；若 \(S\cap\mathfrak p=\varnothing\)，则 \(S^{-1}Q\) 为 \(S^{-1}\mathfrak p\) 准素理想，且其在 \(R\) 内的收缩是 \(Q\)。反之，\(S^{-1}R\) 中任一准素理想的收缩也是准素理想。
\end{proposition}
\begin{proof}
第一种情形取 \(s\in S\cap\mathfrak p\)，则某个 \(s^n\in Q\)，局部化后 \(Q\) 包含单位。第二种情形中，若 \(a/1\in S^{-1}Q\)，则某个 \(s\in S\) 满足 \(sa\in Q\)；因为 \(s\notin\mathfrak p\)，\cref{b3:prop:primary-radical-colon}给出 \(a\in Q\)，这证明收缩断言，也说明扩张仍是真理想。

若 \((a/s)(b/t)\in S^{-1}Q\)、\(b/t\notin S^{-1}Q\)，收缩断言给出 \(ab\in Q\)、\(b\notin Q\)。于是 \(a^n\in Q\)，进而 \((a/s)^n\in S^{-1}Q\)。又由已证的收缩性质，\((a/s)^n\in S^{-1}Q\) 当且仅当 \(a^n\in Q\)。某个这样的正整数 \(n\) 存在，恰等价于 \(a\in\sqrt Q=\mathfrak p\)，所以局部化后的根为 \(S^{-1}\mathfrak p\)。最后，对于收缩的理想，把 \(ab\)、\(b\) 送入局部化环并应用准素条件，所获 \((a/1)^n\) 属于该理想就意味着 \(a^n\) 属于收缩。
\end{proof}

\begin{definition}\label{b3:def:symbolic-power}
设 \(R\) 为交换含幺环，\(\mathfrak p\) 为 \(R\) 的素理想，\(n\ge1\) 为整数。定义
\[
\mathfrak p^{(n)}\coloneqq\mathfrak p^nR_{\mathfrak p}\cap R
\]
并称它为 \(\mathfrak p\) 的第 \(n\) 个\term[fuhao-mi]{符号幂}[symbolic power]；交表示沿 \(R\rightarrow R_{\mathfrak p}\) 的收缩，不要求该映射单射。
\symidx{\(\mathfrak p^{(n)}\)：素理想的符号幂}
\end{definition}

\(\mathfrak pR_{\mathfrak p}\) 是极大理想，所以由\cref{b3:prop:maximal-radical-primary}和\cref{b3:prop:primary-localization}，\(\mathfrak p^{(n)}\) 总是 \(\mathfrak p\) 准素理想，且包含普通幂 \(\mathfrak p^n\)。在\cref{b3:exm:prime-square-not-primary}中，\(\bar y\) 在 \(R_{\mathfrak p}\) 中可逆，故 \(\bar x=\bar z^2/\bar y\in\mathfrak p^2R_{\mathfrak p}\)，于是 \(\bar x\in\mathfrak p^{(2)}\setminus\mathfrak p^2\)。符号幂正是先在目标素理想处集中研究，再把结果收回原环所得的准素对象。

% !TeX root = ../../Book3.tex
% 纲要标题：### 4.3 准素分解定理
% 纲要位置：工作记录/计划纲要.md，第 1919 行。
% 纲要细目（写作范围）：
% * Noether 环中的存在性
% * 极小分解、孤立分量和嵌入分量
% * 哪些部分唯一、哪些部分不唯一

\section{准素分解定理}
\label{b3:sec:0403}

整数的素因数分解给出 \((12)=(4)\cap(3)\)。一般 Noether 环的理想也能分解为有限个准素理想的交，但“各因子都唯一”不再成立。极小分解中的根素理想集合由原理想决定；其中对应极小素理想的准素分量也唯一，嵌入准素分量却可能改变。

\subsection{Noether 环中的准素分解存在性}
\label{b3:subsec:0403:01}

设 \(R\) 为交换含幺环，\(I\subsetneq R\) 为真理想。若对 \(R\) 的任意理想 \(J,K\)，等式 \(I=J\cap K\) 总使 \(I=J\) 或 \(I=K\)，则称 \(I\) 为\term[buke-yue-lixiang]{不可约理想}[irreducible ideal]。这里讨论的是理想作为交的不可约性，不是元素的不可约性，也不是模的单性。

\begin{lemma}\label{b3:lem:irreducible-ideal-primary}
设 \(R\) 为交换含幺 Noether 环。则 \(R\) 的每个真理想是有限个不可约理想的交，而每个不可约理想都是准素理想。
\end{lemma}
\begin{proof}
若有真理想不能写成有限个不可约理想的交，在这些理想中取极大者 \(I\)。它不可能不可约，故 \(I=J\cap K\)，其中 \(J,K\) 都严格包含 \(I\)。极大性保证两者都已有这样的有限分解；将两个分解合并就得到 \(I\) 的分解，矛盾。

设 \(Q\) 不可约且 \(ab\in Q\)、\(b\notin Q\)。升链
\[
(Q:a)\subseteq(Q:a^2)\subseteq\cdots
\]
稳定；选 \(n\ge1\) 使 \((Q:a^n)=(Q:a^{n+1})\)。我们证明
\[
Q=\bigl(Q+(a^n)\bigr)\cap\bigl(Q+(b)\bigr).
\]
只需检查反向包含。若 \(w=q+ra^n=q'+sb\)，其中 \(q,q'\in Q\)，则乘以 \(a\) 后第二式给出 \(aw\in Q\)，于是第一式给出 \(ra^{n+1}\in Q\)。由理想商稳定性，\(ra^n\in Q\)，故 \(w\in Q\)。因为 \(b\notin Q\)，不可约性迫使 \(Q+(a^n)=Q\)，即 \(a^n\in Q\)。这正是准素条件。
\end{proof}

\begin{theorem}[准素分解存在性]\label{b3:thm:primary-decomposition}
设 \(R\) 为交换含幺 Noether 环，\(I\subsetneq R\) 为真理想。则 \(I\) 能写成
\begin{equation}\label{b3:eq:primary-decomposition}
I=Q_1\cap\cdots\cap Q_r,
\end{equation}
其中各 \(Q_i\) 为准素理想。可以进一步要求没有任何 \(Q_i\) 能删去，且各根 \(\mathfrak p_i=\sqrt{Q_i}\) 两两不同。
\end{theorem}
\begin{proof}
\cref{b3:lem:irreducible-ideal-primary}给出一个有限准素分解。把根相同的因子合并为一个交，由\cref{b3:prop:primary-radical-colon}它仍为同根准素理想。再删去冗余因子，即得附加条件。
\end{proof}

单项式理想的成员关系可以逐项判断。设 \(k\) 为域，在 \(k[x_1,\ldots,x_n]\) 中，由单项式 \(m_1,\ldots,m_s\) 生成的理想，作为 \(k\) 向量空间，恰由那些被某个 \(m_i\) 整除的单项式生成。事实上，把各多项式系数展开，就得到这一方向的包含；反向，每个这样的单项式本来就是某个 \(m_i\) 的倍数。由于不同单项式线性无关，一个多项式属于此理想，当且仅当它的每个非零单项都属于此理想。因此两个单项式理想取交时可以逐项检查，不需要预先运行 Gröbner 基算法。

\ProseParagraphBreak
准素理想不一定不可约。例如，设 \(k\) 为域，在 \(k[x,y]\) 内有
\[
(x,y)^2=(x^2,y)\cap(x,y^2).
\]
等式可逐个比较单项式：同时属于右端的单项式恰被 \(x^2,xy,y^2\) 之一整除。两边的理想都是 \((x,y)\) 准素理想，右边两个因子却都严格大于左边。因此不可约分解只是证明存在性的一种工具，并不是最后要保存的标准形式。

\subsection{极小分解、孤立分量与嵌入分量}
\label{b3:subsec:0403:02}

设 \(R\) 为交换含幺环，\(I\subsetneq R\) 为真理想。若 \(I=Q_1\cap\cdots\cap Q_r\) 的各因子均为准素理想、各根 \(\sqrt{Q_i}\) 两两不同，且删去任一因子都会改变交，则称它为\term[jixiao-zhunsu-fenjie]{极小准素分解}[minimal primary decomposition]，其中各个 \(Q_i\) 称为\term[zhunsu-fenliang]{准素分量}[primary component]。“极小”不声称每个准素因子在包含关系下极小。

\begin{theorem}\label{b3:thm:primary-associated-primes}
设 \(R\) 为交换含幺 Noether 环，\(I\subsetneq R\) 为真理想。若 \(I=\bigcap_{i=1}^rQ_i\) 是 \(R\) 中的极小准素分解，其中 \(r\ge1\)，各 \(Q_i\) 为准素理想，\(\sqrt{Q_i}=\mathfrak p_i\)，则
\[
\operatorname{Ass}_R(R/I)=\{\mathfrak p_1,\ldots,\mathfrak p_r\}.
\]
特别地，分解中出现哪些素理想与分解的选择无关。
\end{theorem}
\begin{proof}
对角映射给出单射
\[
R/I\longrightarrow\bigoplus_{i=1}^rR/Q_i,\qquad
\bar a\longmapsto(a+Q_1,\ldots,a+Q_r).
\]
由\cref{b3:prop:associated-exact-sequence}和\cref{b3:prop:primary-associated-singleton}，左端的关联素理想都在 \(\{\mathfrak p_i\}\) 中。

固定 \(i\)，令 \(J_i=\bigcap_{j\ne i}Q_j\)，在 \(r=1\) 时约定 \(J_i=R\)。无冗余保证 \(J_i/I\ne0\)，而映射 \(J_i/I\rightarrow R/Q_i\) 单射，因其核为 \((J_i\cap Q_i)/I=0\)。由\cref{b3:thm:associated-existence}，这个非零子模有至少一个关联素理想；由于它嵌入 \(R/Q_i\)，该素理想只能是 \(\mathfrak p_i\)。它同时是 \(R/I\) 的子模，故 \(\mathfrak p_i\in\operatorname{Ass}_R(R/I)\)。逐个取 \(i\) 得反向包含。
\end{proof}

按照前面对关联素理想的命名，若 \(\mathfrak p_i\) 为孤立关联素理想，则称相应 \(Q_i\) 为孤立准素分量；若 \(\mathfrak p_i\) 为嵌入关联素理想，则称 \(Q_i\) 为嵌入准素分量。由\cref{b3:prop:minimal-primes-associated}，孤立关联素理想恰是包含 \(I\) 的极小素理想。它们对应 \(V(I)\) 的不可约分量，而嵌入素理想对应的闭集位于某个不可约分量之内，并不另给出一个不可约分量。

\ProseParagraphBreak
例如，设 \(k\) 为域，理想 \(I=(x^2,xy)\subset k[x,y]\) 满足
\begin{equation}\label{b3:eq:embedded-line-decomposition}
I=(x)\cap(x^2,y).
\end{equation}
第二因子因根为极大理想 \((x,y)\) 而准素。要验证交，设 \(xf=x^2a+yb\)，模掉 \(x\) 后有 \(y\bar b=0\) 于 \(k[y]\)，故 \(b=xc\)，于是 \(xf=x^2a+xyc\in I\)。两因子均不可删，根之间的包含关系为
\[
\underbrace{(x)}_{\text{极小关联素理想}}
\subsetneq
\underbrace{(x,y)}_{\text{嵌入关联素理想}}.
\]
因此 \((x)\) 是孤立准素分量，\((x^2,y)\) 是嵌入准素分量；\(V(I)=V(x)\) 只有一个不可约分量，\(V(x,y)\) 是其中的一个闭点。

\subsection{准素分解的唯一性与非唯一性}
\label{b3:subsec:0403:03}

\begin{theorem}\label{b3:thm:primary-isolated-uniqueness}
设 \(R\) 为交换含幺 Noether 环，\(I\subsetneq R\) 为真理想，\(I=\bigcap_{i=1}^rQ_i\) 为极小准素分解，其中 \(r\ge1\)、\(\mathfrak p_i=\sqrt{Q_i}\)。若 \(\mathfrak p_i\) 为 \(\operatorname{Ass}_R(R/I)\) 中的极小元，则
\[
Q_i=IR_{\mathfrak p_i}\cap R.
\]
因此每个孤立准素分量由 \(I\) 与其根唯一决定。
\end{theorem}
\begin{proof}
有限交与局部化交换。若 \(j\ne i\)，由于 \(\mathfrak p_i\) 在所有根中极小且根互异，\(\mathfrak p_j\not\subseteq\mathfrak p_i\)。取 \(a\in\mathfrak p_j\setminus\mathfrak p_i\)，则某个 \(a^n\in Q_j\)，而 \(a\) 在 \(R_{\mathfrak p_i}\) 中成为单位，所以 \(Q_jR_{\mathfrak p_i}=R_{\mathfrak p_i}\)。于是 \(IR_{\mathfrak p_i}=Q_iR_{\mathfrak p_i}\)。\cref{b3:prop:primary-localization}又说明 \(Q_i\) 是其扩张的收缩，得到公式。
\end{proof}

嵌入准素分量一般不唯一。在\cref{b3:eq:embedded-line-decomposition}的同一个理想上，对每个 \(n\ge1\) 都有
\begin{equation}\label{b3:eq:nonunique-embedded-components}
(x^2,xy)=(x)\cap(x^2,xy,y^n).
\end{equation}
右边第二因子的根为 \((x,y)\)，故它准素。由单项式整除判据，第二因子中同时被 \(x\) 整除的单项式恰属于 \((x^2,xy)\)，所以交确实不变。\(x\) 不属于第二因子，\(y^n\) 不属于第一因子，因而各分解都极小；而 \(y^n\) 属于第 \(n\) 个第二因子，却不属于第 \(n+1\) 个，故这些准素分量互不相同。

\ProseParagraphBreak
极小准素分解中出现的根素理想集合唯一，并等于 \(\operatorname{Ass}_R(R/I)\)；孤立准素分量唯一，嵌入准素分量一般不唯一。\cref{b3:tab:primary-uniqueness-data}区分了这些数据。模的对应结论及另一种存在性证明见下一节；文献中的统一表述可对照 \cite[\S6, Theorem 6.8]{Matsumura1989CommutativeRingTheory}。

\begin{table}[htbp]
\centering
\caption[极小准素分解的唯一性数据]{极小准素分解的唯一性数据。设 \(R\) 为交换含幺 Noether 环，\(I\subsetneq R\) 为真理想。}
\label{b3:tab:primary-uniqueness-data}
\begin{tabularx}{\textwidth}{@{}>{\raggedright\arraybackslash}X>{\raggedright\arraybackslash}p{.27\textwidth}@{}}
\toprule
数据 & 是否由 \(I\) 唯一确定\\
\midrule
极小准素分解中的根素理想集合 \(\operatorname{Ass}_R(R/I)\) & 是\\
孤立关联素理想集合 & 是\\
对应于给定孤立关联素理想的准素分量 & 是\\
嵌入关联素理想集合 & 是\\
对应于给定嵌入关联素理想的准素分量 & 一般不唯一\\
\bottomrule
\end{tabularx}
\end{table}

% !TeX root = ../../Book3.tex
% 纲要标题：### 4.4 子模的准素分解与嵌入分量
% 纲要位置：工作记录/计划纲要.md，第 2026 行。
% 纲要细目（写作范围）：
% * 子模的准素分解；商模上的统一幂零作用及有限生成假设的无限直和反例
% * 嵌入关联素理想的显式例子；实际元素零化子及矩阵余核的生成关系
% * 与几何分支及非约化结构的联系；区分多分支、沿分量非约化与特殊点处的嵌入结构
% ---

\section{子模的准素分解与嵌入分量}
\label{b3:sec:0404}

理想分解其实是自由模 \(R\) 中的子模分解。若要研究一个矩阵给出的余核，或研究有限生成模的支集与局部结构，就需要允许一般有限生成模，并把准素条件写成商模上标量作用的性质。

\subsection{子模的准素分解}
\label{b3:subsec:0404:01}

\begin{definition}\label{b3:def:primary-submodule}
设 \(R\) 为交换含幺环，\(M\) 为 \(R\) 模，\(N\subsetneq M\) 为真子模。若对任意 \(a\in R\)、\(m\in M\) 都有
\[
am\in N,\quad m\notin N
\quad\Longrightarrow\quad a^nM\subseteq N
\ \text{对某个}\;n\ge1.
\]
则称 \(N\) 为\term[zhunsu-zimo]{准素子模}[primary submodule]。若 \(N\) 准素，且 \(\sqrt{\operatorname{Ann}_R(M/N)}\) 等于素理想 \(\mathfrak p\)，则称 \(N\) 为 \(\mathfrak p\) 准素子模。整个条件仅依赖商模 \(M/N\)，但“\(N\) 准素”总须指明它位于哪个模中。
\end{definition}

令 \(T=M/N\)。条件 \(am\in N\)、\(m\notin N\) 表示 \(a\) 零化了 \(T\) 中的某个非零元素，因而是 \(T\) 上的零因子；结论 \(a^nM\subseteq N\) 则表示 \(a^nT=0\)。所以准素性要求每个零因子的某次幂零化整个商模。指数可以随标量 \(a\) 改变，但选定 \(a\) 后，必须有同一个指数适用于 \(T\) 的所有元素。仅要求眼前的 \(m\) 被 \(a^n\) 零化并不能表达这一点，因为 \(am\in N\) 已经使它成立。若只知道对每个 \(t\in T\) 分别存在指数使 \(a^{n(t)}t=0\)，在 \(T\) 有限生成时，可对有限组生成元取指数的最大值，得到 \(a^nT=0\)。当 \(T=R/Q\) 时，\(a^nT=0\) 等价于 \(a^n\in Q\)，故此定义与准素理想的定义一致。

\begin{proposition}\label{b3:prop:primary-submodule-associated}
设 \(R\) 为交换含幺 Noether 环，\(M\) 为有限生成 \(R\) 模，\(N\subsetneq M\) 为真子模。\(N\) 准素，当且仅当 \(\operatorname{Ass}_R(M/N)\) 只有一个元素；此元素就是 \(\sqrt{\operatorname{Ann}_R(M/N)}\)，因而该根是素理想。\(M\) 的有限生成假设一般不能删去。
\end{proposition}
\begin{proof}
记 \(T=M/N\) 与 \(J=\operatorname{Ann}_R(T)\)。若 \(N\) 准素，取 \(\mathfrak p=\operatorname{Ann}_R(t)\in\operatorname{Ass}_R(T)\)。因 \(J\subseteq\mathfrak p\)，有 \(\sqrt J\subseteq\mathfrak p\)。另一方面，每个 \(a\in\mathfrak p\) 都零化非零元 \(t\)，准素条件给出 \(a^nT=0\)，即 \(a\in\sqrt J\)。故所有关联素理想都等于 \(\sqrt J\)，且存在性保证集合非空。

反向若 \(\operatorname{Ass}_R(T)=\{\mathfrak p\}\)，这里 \(T\) 的有限生成性保证 \(\operatorname{Supp}_R(T)=V(J)\)。由\cref{b3:prop:minimal-primes-associated}，\(V(J)\) 的极小素理想只有 \(\mathfrak p\)，且其中每个素理想都包含 \(\mathfrak p\)。于是包含 \(J\) 的全部素理想之交就是 \(\mathfrak p\)，即 \(\sqrt J=\mathfrak p\)。若 \(at=0\)、\(t\ne0\)，由\cref{b3:thm:associated-zero-divisors}得 \(a\in\mathfrak p=\sqrt J\)，故某个 \(a^n\in J\)，这就是 \(a^nM\subseteq N\)。

为说明有限生成假设的作用，取域 \(k\)，令 \(R=k[x]\) 与
\[
T=\bigoplus_{n\ge1}R/(x^n).
\]
每个元素只有有限个非零坐标，因此被 \(x\) 的某次幂零化。有限个元素总共只涉及有限个直和分量，不能生成全部 \(T\)，故 \(T\) 不是有限生成模。若非零元素 \(t\) 的零化子是素理想 \(\mathfrak q\)，由 \(x^nt=0\) 得 \(x\in\mathfrak q\)；因 \((x)\) 是极大理想，只能有 \(\mathfrak q=(x)\)。而第一坐标中 \(1+(x)\) 的零化子恰为 \((x)\)，故 \(\operatorname{Ass}_R(T)=\{(x)\}\)。\(x\) 零化这个非零元素，所以是 \(T\) 上的零因子；但对任意 \(n\ge1\)，第 \(n+1\) 个坐标中的 \(1+(x^{n+1})\) 不被 \(x^n\) 零化。因而不存在统一指数使 \(x^nT=0\)，\(0\subset T\) 不是准素子模。逐元素的指数在这里无法统一，且 \(\operatorname{Ann}_R(T)=\bigcap_{n\ge1}(x^n)=0\)，其根也不同于唯一的关联素理想 \((x)\)。
\end{proof}

\begin{theorem}\label{b3:thm:module-primary-decomposition}
设 \(R\) 为交换含幺 Noether 环，\(M\) 为有限生成 \(R\) 模。则 \(M\) 的每个真子模 \(N\) 都有极小准素分解
\[
N=N_1\cap\cdots\cap N_r,
\qquad \operatorname{Ass}_R(M/N_i)=\{\mathfrak p_i\},
\]
其中各 \(\mathfrak p_i\) 互异，且无因子可删。并且
\[
\operatorname{Ass}_R(M/N)=\{\mathfrak p_1,\ldots,\mathfrak p_r\}.
\]
若 \(\mathfrak p_i\) 为该集合的极小元，记 \(\lambda_i:M\rightarrow M_{\mathfrak p_i}\) 为局部化映射，则
\[
N_i=\lambda_i^{-1}(N_{\mathfrak p_i}).
\]
\end{theorem}
\begin{proof}
若真子模 \(N\subsetneq M\) 不能写成两个严格更大的子模的交，则称 \(N\) 在 \(M\) 中不可约。因为 \(M\) 是 Noether 模，\cref{b3:lem:irreducible-ideal-primary}第一段的极大反例证明逐字适用于子模，故 \(N\) 是有限个不可约子模的交。这个有限分解提供了候选因子，还需证明这些因子准素，并将其整理为定理要求的极小准素分解。

不可约子模必准素。否则由\cref{b3:prop:primary-submodule-associated}，非零模 \(M/N\) 至少有两个不同关联素理想 \(\mathfrak p,\mathfrak q\)。分别取子模 \(L\cong R/\mathfrak p\)、\(K\cong R/\mathfrak q\)。若 \(0\ne z\in L\cap K\)，在 \(L\cong R/\mathfrak p\) 中将 \(z\) 对应于非零剩余类 \(\bar b\)。因 \(R/\mathfrak p\) 为整环，\(az=0\) 等价于 \(\bar a\bar b=0\)，也就等价于 \(a\in\mathfrak p\)。所以 \(\operatorname{Ann}_R(z)=\mathfrak p\)；同理，由 \(z\in K\) 又得 \(\operatorname{Ann}_R(z)=\mathfrak q\)，矛盾。所以 \(L\cap K=0\)，两个非零子模在 \(M\) 中的原像便把 \(N\) 写成两个严格更大子模的交，与不可约性矛盾。

为合并根相同的因子，考虑任意非空有限族同根准素子模的交。其商模单射到各商模的直和，由\cref{b3:prop:associated-exact-sequence}，关联素理想只能属于这个单元素集合；商模非零，由\cref{b3:thm:associated-existence}恰有该关联素理想，再由\cref{b3:prop:primary-submodule-associated}得到准素性。于是合并同根因子并删除冗余后，就得到根两两不同、无因子可删的分解。

这些根确实都是 \(M/N\) 的关联素理想。对角单射 \(M/N\hookrightarrow\bigoplus_iM/N_i\) 先给出 \(\operatorname{Ass}_R(M/N)\subseteq\{\mathfrak p_1,\ldots,\mathfrak p_r\}\)。反向，去掉第 \(i\) 个因子会使交严格增大，故 \((\bigcap_{j\ne i}N_j)/N\) 非零；当 \(r=1\) 时，此处空族的交取为 \(M\)。它是 \(M/N\) 的子模，而映入 \(M/N_i\) 的核是 \((\bigcap_jN_j)/N=0\)，所以也嵌入 \(M/N_i\)。它的关联素理想非空且只能为 \(\mathfrak p_i\)，故每个 \(\mathfrak p_i\) 都在 \(\operatorname{Ass}_R(M/N)\) 中出现。

若 \(\mathfrak p_i\) 极小，在 \(\mathfrak p_i\) 处局部化可以使其余分量的商模消失。对 \(j\ne i\)，极小性与根互异保证 \(\mathfrak p_j\not\subseteq\mathfrak p_i\)，故可取 \(a\in\mathfrak p_j\setminus\mathfrak p_i\)。因为 \(\mathfrak p_j=\sqrt{\operatorname{Ann}_R(M/N_j)}\)，某个 \(a^n\) 零化 \(M/N_j\)。局部化后 \(a\) 是单位，故 \((N_j)_{\mathfrak p_i}=M_{\mathfrak p_i}\)。局部化与有限交相容，于是 \(N_{\mathfrak p_i}=(N_i)_{\mathfrak p_i}\)。再将这个局部子模收缩：若 \(m/1\in(N_i)_{\mathfrak p_i}\)，则某个 \(s\notin\mathfrak p_i\) 满足 \(sm\in N_i\)。若 \(m\notin N_i\)，准素性将给出 \(s\in\sqrt{\operatorname{Ann}_R(M/N_i)}=\mathfrak p_i\)，矛盾。所以收缩恰为 \(N_i\)，公式得证。
\end{proof}

\subsection{嵌入关联素理想的显式例子}
\label{b3:subsec:0404:02}

取域 \(k\)，回到 \(R=k[x,y]\)、\(I=(x^2,xy)\)。商模 \(R/I\) 中有
\[
\operatorname{Ann}_R(\bar x)=(x,y),\qquad
\operatorname{Ann}_R(\bar y)=(x).
\]
第一式因为 \(f\bar x=f(0,0)\bar x\)；第二式可将 \(fy\in(x^2,xy)\) 模掉 \(x\)，由 \(k[y]\) 是整环得到 \(f\in(x)\)。这不仅列出了素理想，还给出了定义它们的实际元素。它们正是\cref{b3:thm:primary-associated-primes}对\cref{b3:eq:embedded-line-decomposition}所给出的两个关联素理想。

\ProseParagraphBreak
一个矩阵也可以同时包含孤立与嵌入成分。令
\[
\Phi:R^3\longrightarrow R^2,\qquad
[\Phi]=\begin{pmatrix}x&y&0\\0&0&x\end{pmatrix},
\qquad T=\operatorname{coker}\Phi.
\]
设 \(e_1,e_2\) 为 \(R^2\) 的标准基，\(\bar e_1,\bar e_2\) 为它们在余核中的像。矩阵的三列给出关系
\[
x\bar e_1=0,\qquad y\bar e_1=0,\qquad x\bar e_2=0.
\]
像子模为 \(\operatorname{im}\Phi=(x,y)e_1\oplus(x)e_2\)，两组关系没有混合两个坐标，故 \(R\bar e_1\cong R/(x,y)\)、\(R\bar e_2\cong R/(x)\)，并且
\[
T\cong R/(x,y)\oplus R/(x).
\]
因此 \(\operatorname{Ass}_R(T)=\bigl\{(x,y),(x)\bigr\}\)。\(N=\operatorname{im}\Phi\) 的极小准素分解为
\[
N=\bigl(Re_1\oplus(x)e_2\bigr)
   \cap\bigl((x,y)e_1\oplus Re_2\bigr).
\]
第一个商模为 \(R/(x)\)，第二个为 \(R/(x,y)\)，故两因子分别是 \((x)\) 与 \((x,y)\) 准素分量。在 \((x)\) 处局部化时 \(y\) 成为单位，第一坐标的商消失，第二坐标留下 \(k(y)\)。这就是孤立准素分量的收缩公式在一个展示矩阵上的具体表现。

\subsection{几何分支与非约化结构}
\label{b3:subsec:0404:03}

设 \(R\) 为交换含幺 Noether 环，\(I\subsetneq R\) 为真理想。闭集 \(V(I)=V(\sqrt I)\) 的不可约分量由包含 \(I\) 的极小素理想决定，其中仍包含这些极小点的所有特殊化点，包括嵌入关联素理想对应的点。底拓扑却不能辨认一个非极小点是否为嵌入关联点：将 \(I\) 换为 \(\sqrt I\) 不改变这个闭集，却会除去商环中的幂零元，并可能改变关联素理想集合。以下直接在素谱中解释这种区别，无须预设基域代数闭。

\begin{proposition}\label{b3:prop:reduced-ring-no-embedded-primes}
设 \(R\) 为交换含幺约化环。则
\[
\operatorname{Ass}_R(R)\subseteq\min\operatorname{Spec}R.
\]
因此任意约化环都没有嵌入关联素理想。若 \(R\) 还为 Noether 环，则上述包含为等号。
\end{proposition}
\begin{proof}
若 \(R=0\)，两侧均为空集。以下设 \(R\ne0\)。取 \(\mathfrak p=\operatorname{Ann}_R(x)\in\operatorname{Ass}_R(R)\)，其中 \(x\ne0\)。若 \(\mathfrak p\) 不是极小素理想，则存在素理想 \(\mathfrak q\subsetneq\mathfrak p\)。取 \(a\in\mathfrak p\setminus\mathfrak q\)，由 \(ax=0\in\mathfrak q\) 和 \(a\notin\mathfrak q\)，素性给出 \(x\in\mathfrak q\subseteq\mathfrak p=\operatorname{Ann}_R(x)\)。于是 \(x^2=0\)，约化性迫使 \(x=0\)，矛盾。这证明了包含式。

若 \(R\) 为 Noether 环，对有限生成模 \(R\) 使用\cref{b3:prop:minimal-primes-associated}；因 \(\operatorname{Supp}_R(R)=\operatorname{Spec}R\)，每个极小素理想都属于 \(\operatorname{Ass}_R(R)\)，从而得到等号。
\end{proof}

取域 \(k\)，先看 \(C=k[x,y]/(xy)\)。因为 \((xy)=(x)\cap(y)\)，它是约化环，有两个极小素理想 \((\bar x)\)、\((\bar y)\)，对应两条坐标轴这两个不可约分量。由刚才的命题，它的关联素理想也恰为这两个极小素理想，没有嵌入关联素理想。这里的两个分支来自两个极小素理想，并不来自幂零结构。

\ProseParagraphBreak
即使在 Noether 环中，没有嵌入关联素理想也不保证约化。环 \(B=k[x,y]/(x^2)\) 非约化，却只有孤立关联素理想 \((\bar x)\)。每个元素唯一写成 \(a(y)+\bar x\cdot b(y)\)，其中 \(a,b\in k[y]\)。对 \(c,d\in k[y]\)，有
\[
(a+\bar x\cdot b)(c+\bar x\cdot d)=ac+\bar x\cdot(ad+bc).
\]
若 \(a\ne0\) 且乘积为零，比较两个 \(k[y]\) 分量得 \(ac=0\)、\(ad+bc=0\)。因为 \(k[y]\) 是整环，先得 \(c=0\)，再得 \(d=0\)；故 \(a+\bar x\cdot b\) 的乘法是单射。零因子因此恰为 \((\bar x)\) 中的元素，且都平方为零，说明 \((x^2)\subset k[x,y]\) 为 \((x)\) 准素理想。

\ProseParagraphBreak
在 \(B\) 中，幂零理想 \((\bar x)\) 作为 \(B\) 模同构于 \(B/(\bar x)\cong k[y]\)，其支集是整条直线 \(V(\bar x)\)，因而在这条线的每个点处局部化都不会消失。再看 \(A=k[x,y]/(x^2,xy)\)，其幂零理想 \((\bar x)\) 被 \((\bar x,\bar y)\) 零化，且作为 \(k\) 向量空间只有一维，所以支集只有原点；在开集 \(D(\bar y)\) 上它消失。上一小节计算的 \(\operatorname{Ann}_{k[x,y]}(\bar x)=(x,y)\) 正是这个原点处嵌入关联素理想的来源，它没有增加新的不可约分量。

\ProseParagraphBreak
\(A\) 与 \(B\) 的约化商都同构于 \(k[y]\)，底空间只有一条不可约直线，幂零理想的支集却不同。\(B\) 没有嵌入关联素理想，仍然非约化；\(A\) 在原点有嵌入关联素理想，仍然只有一个不可约分量。相较之下，约化环 \(C\) 有两个几何分支，却没有幂零结构。


\begin{chaptersummary}
关联素理想由循环子模 \(R/\mathfrak p\) 的实际嵌入定义。Noether 环保证非零模有这样的子模，有限生成模的升链条件进一步给出有限素理想滤过，因而关联素理想只有有限个。它们的并是全部零因子，它们的极小元与支集的极小元一致；局部化保留恰好那些不遇到分母集的关联素理想。

\ProseParagraphBreak
在 Noether 环中，理想及有限生成模中的子模可以写成有限个准素分量的交。极小分解中出现的根素理想集合由原商模唯一确定；孤立准素分量还可以用局部化及收缩唯一恢复，嵌入准素分量则一般不唯一。任意约化交换环的关联素理想都是极小素理想，Noether 情形下两者的集合相等；反过来，没有嵌入关联素理想的环仍可能带有幂零元。
\end{chaptersummary}

\BookExercises
\begin{exercise}\label{b3:ex:04-integer-module}
求 \(M=\mathbb Z\oplus\mathbb Z/12\mathbb Z\) 的关联素理想与零因子，并为每个关联素理想给出一个零化子恰为它的元素。
\end{exercise}
\begin{solution}
自由因子的非零元素零化子为 \((0)\)。第二因子内，\(\bar a\) 的零化子由 \(12/\gcd(12,a)\) 生成，成为素理想时只能是 \((2)\) 或 \((3)\)。因此关联素理想为 \((0),(2),(3)\)，可分别取 \((1,\bar0),(0,\bar6),(0,\bar4)\) 为见证。由\cref{b3:thm:associated-zero-divisors}，零因子为 \((0)\cup(2)\cup(3)=(2)\cup(3)\)，即 \(2\mathbb Z\cup3\mathbb Z\)。关联素理想 \((0)\) 包含于另两个理想，没有向并集增加元素；若只保留模 \(\mathbb Z/12\mathbb Z\)，关联集合便少了 \((0)\)，零因子集合却相同。因此仅由零因子的并集不能恢复全部关联素理想。
\end{solution}

\begin{exercise}\label{b3:ex:04-primary-polynomial-module}
设 \(k\) 为域。在 \(R=k[x,y]\) 中，直接计算 \(R/(x^3)\) 的零因子及关联素理想，并求元素 \(\bar x^2\) 的零化子。
\end{exercise}
\begin{solution}
每个类唯一写成 \(a_0(y)+a_1(y)\bar x+a_2(y)\bar x^2\)。若 \(a_0\ne0\)，将其与另一个这样的式子相乘，逐次比较 \(1,\bar x,\bar x^2\) 的系数，利用 \(k[y]\) 为整环可知乘法单射。若 \(a_0=0\)，该元素是幂零元，且会零化非零元 \(\bar x^2\)。所以零因子来自 \((x)\)，\((x^3)\) 为 \((x)\) 准素理想，关联素理想集合为 \(\bigl\{(x)\bigr\}\)。计算 \(f\bar x^2=a_0(y)\bar x^2\) 给出零化子 \((x)\)。
\end{solution}

\begin{exercise}\label{b3:ex:04-exact-quotient-warning}
设 \(k\) 为域，\(R=k[t]\)，令 \(M=R/(t^2)\oplus R\)、\(N=R/(t^2)\oplus t(t-1)R\subset M\)、\(P=M/N\)。计算这三个模的关联素理想，并验证\cref{b3:prop:associated-exact-sequence}中的两个包含。说明第二个包含可以严格，且商模既可以产生新的关联素理想，也可以失去原有的关联素理想。
\end{exercise}
\begin{solution}
因 \(t(t-1)R\cong R\)，有限直和公式给出
\[
\operatorname{Ass}_R(N)=\operatorname{Ass}_R(M)=\{(0),(t)\}.
\]
商模为 \(P\cong R/\bigl(t(t-1)\bigr)\)。理想 \((t)\)、\((t-1)\) 互素，中国剩余定理给出 \(P\cong R/(t)\oplus R/(t-1)\)，所以
\[
\operatorname{Ass}_R(P)=\{(t),(t-1)\}.
\]
第一个包含在此取等号，而第二个为
\(\{(0),(t)\}\subsetneq\{(0),(t),(t-1)\}\)。商模增加了 \((t-1)\)，同时失去 \((0)\)；两个方向都不能从原模的关联集合直接推定。
\end{solution}

\begin{exercise}\label{b3:ex:04-infinite-associated}
求 \(\mathbb Q/\mathbb Z\) 作为 \(\mathbb Z\) 模的关联素理想。说明它的例子不违反关联素理想有限性定理。
\end{exercise}
\begin{solution}
非零元 \(a/n+\mathbb Z\) 在 \(\gcd(a,n)=1\) 时零化子为 \((n)\)，故素零化子恰为所有素数 \(p\) 对应的 \((p)\)，由 \(1/p+\mathbb Z\) 实现。不存在零化子为零的元素，因为每个元素都有有限加法阶。所得集合无限，而 \(\mathbb Q/\mathbb Z\) 不有限生成：有限多个有理数的分母有公倍数 \(d\)，它们生成的模被 \(d\) 零化，却不能包含所有 \(1/p+\mathbb Z\)。有限性定理的模假设不满足。
\end{solution}

\begin{exercise}\label{b3:ex:04-localization-primes}
设 \(k\) 为域，令 \(R=k[x,y]\)、\(M=R/(x)\oplus R/(x,y)\)，取 \(\mathfrak p=(x)\)、\(\mathfrak m_0=(x,y)\)、\(\mathfrak m_1=(x,y-1)\)。计算这三个素理想处的局部模、关联素理想，以及 \(M_{\mathfrak q}/\mathfrak qM_{\mathfrak q}\) 在剩余域 \(\kappa(\mathfrak q)=R_{\mathfrak q}/\mathfrak qR_{\mathfrak q}\) 上的维数，其中 \(\mathfrak q\) 分别取上述三个理想。说明局部环的极大理想何时是该模的关联素理想。
\end{exercise}
\begin{solution}
原关联素理想为 \((x),(x,y)\)。在 \(\mathfrak p\) 处，\(y\) 成为单位，第二因子消失，第一因子将 \(k[y]\) 中所有非零元素倒置，故 \(M_{\mathfrak p}\cong k(y)\)。关联集合为 \(\{(x)R_{\mathfrak p}\}\)，该理想就是局部环的极大理想；模已被它零化，剩余向量空间在 \(k(y)\) 上维数为一。

在 \(\mathfrak m_0\) 处，两因子都保留，得到 \(M_{\mathfrak m_0}\cong k[y]_{(y)}\oplus k\)。关联集合为 \(\{(x)R_{\mathfrak m_0},\mathfrak m_0R_{\mathfrak m_0}\}\)，其中极大理想来自第二因子。对两因子分别模掉 \(\mathfrak m_0\)，各得到一个 \(k\)，故剩余向量空间维数为二。

在 \(\mathfrak m_1\) 处，\(y\notin\mathfrak m_1\)，第二因子消失，第一因子为 \(k[y]_{(y-1)}\)。关联集合为 \(\{(x)R_{\mathfrak m_1}\}\)，而 \(y-1\) 不在 \((x)R_{\mathfrak m_1}\) 中，所以这个关联素理想严格小于局部环的极大理想。剩余向量空间仍是一个 \(k\)，维数为一；维数非零并不说明极大理想一定关联。
\end{solution}

\begin{exercise}\label{b3:ex:04-prime-filtration-explicit}
设 \(k\) 为域。对 \(R=k[x,y]\)、\(M=R/(x^2,xy)\)，构造长度为二的素理想滤过，写出两个因子的具体同构。
\end{exercise}
\begin{solution}
取 \(M_1=R\bar x\)。因 \(\operatorname{Ann}_R(\bar x)=(x,y)\)，映射 \(R/(x,y)\rightarrow M_1\)、\(\bar a\mapsto a\bar x\) 同构。商掉 \(M_1\) 相当于再加关系 \(x=0\)，故 \(M/M_1\cong R/(x)\)。链 \(0\subsetneq M_1\subsetneq M\) 即所求；两个商都非零，且都是相应素商环。
\end{solution}

\begin{exercise}\label{b3:ex:04-primary-colon-basis}
设 \(k\) 为域，在 \(R=k[x,y]\) 中令 \(Q=(x^2,y^3)\)。证明 \(Q\) 准素，写出 \(R/Q\) 的一组 \(k\) 基，并计算 \((Q:x)\)、\((Q:y^2)\)。
\end{exercise}
\begin{solution}
\(\sqrt Q=(x,y)\) 极大，所以由\cref{b3:prop:maximal-radical-primary}，\(Q\) 为 \((x,y)\) 准素理想。标准单项式为 \(1,y,y^2,x,xy,xy^2\)，它们生成商空间且因原理想为单项式理想而线性无关。逐个判断乘积所属关系得 \((Q:x)=(x,y^3)\)、\((Q:y^2)=(x^2,y)\)。两理想的根仍为 \((x,y)\)，符合准素理想商命题。
\end{solution}

\begin{exercise}\label{b3:ex:04-prime-radical-insufficient}
设 \(k\) 为域，令 \(R=k[x,y,z]\)、\(J=(x^2,xz)\)。求其根，并用两个明确元素证明 \(J\) 不准素。指出这说明关于根的哪一个逆命题错误。
\end{exercise}
\begin{solution}
\(x^2\in J\subseteq(x)\)，故 \(\sqrt J=(x)\)，为素理想。但 \(zx\in J\)、\(x\notin J\)，而 \(z^n\notin(x)\supseteq J\) 对所有 \(n\ge1\) 成立，所以 \(J\) 不准素。因而“准素理想的根是素理想”的逆命题错误；极大根准素命题的“极大”不能换成“素”。
\end{solution}

\begin{exercise}\label{b3:ex:04-primary-localization-calculation}
设 \(k\) 为域，在 \(R=k[x,y]\) 中令 \(Q=(x^2,y^3)\)。分别倒置 \(1+x\) 与 \(x\)，计算 \(Q\) 的扩张及收缩；说明两种答案的差别。
\end{exercise}
\begin{solution}
\(1+x\notin(x,y)=\sqrt Q\)，因此在 \(R_{1+x}\) 中，扩张是由 \(x^2,y^3\) 生成的真准素理想，收缩仍为 \(Q\)。商环中 \((1+x)(1-x)=1\)，所以 \(1+x\) 原已可逆，得到 \(R_{1+x}/QR_{1+x}\cong R/Q\)。不变的是这个商环；环境环 \(R_{1+x}\) 则严格包含 \(R\)，因为 \(1/(1+x)\) 不是多项式。倒置 \(x\) 时 \(x^2\in Q\) 成为单位，扩张是全环，收缩是 \(R\)。区别恰在分母集是否与根相交。
\end{solution}

\begin{exercise}\label{b3:ex:04-symbolic-square-cone}
设 \(k\) 为域，\(R=k[x,y,z]/(xy-z^2)\)，\(\mathfrak p=(\bar x,\bar z)\)。证明符号平方 \(\mathfrak p^{(2)}=(\bar x)\)。
\end{exercise}
\begin{solution}
商环 \(R/(\bar x)\cong k[y,z]/(z^2)\) 中，每个元素为 \(a(y)+z\cdot b(y)\)。若 \(a\ne0\)，比较乘积的两个系数可见其为非零因子；若 \(a=0\)，它平方为零。故 \((\bar x)\) 为 \(\mathfrak p\) 准素理想。局部化于 \(\mathfrak p\) 后 \(\bar y\) 可逆，\(\bar x=\bar z^2/\bar y\)，且 \(\mathfrak pR_{\mathfrak p}=(\bar z)\)，于是 \(\mathfrak p^2R_{\mathfrak p}=(\bar x)R_{\mathfrak p}\)。利用准素理想的收缩性质得到 \(\mathfrak p^{(2)}=(\bar x)\)。这也再次显示 \(\mathfrak p^2\subsetneq\mathfrak p^{(2)}\)。
\end{solution}

\begin{exercise}\label{b3:ex:04-two-isolated-components}
设 \(k\) 为域。在 \(R=k[x,y,z]\) 中分解 \(I=(xy,xz)\)，求 \(\operatorname{Ass}_R(R/I)\)，并解释其准素分量为何全部唯一。
\end{exercise}
\begin{solution}
\(I=(x)\cap(y,z)\)：若 \(xf\in(y,z)\)，模掉 \((y,z)\) 后在 \(k[x]\) 中有 \(x\bar f=0\)，故 \(f\in(y,z)\)。两个因子都是素理想，互不包含，且 \(x\)、\(y\) 分别见证无冗余。因此关联素理想恰为 \((x),(y,z)\)，两者都极小，准素分量由\cref{b3:thm:primary-isolated-uniqueness}唯一决定。
\end{solution}

\begin{exercise}\label{b3:ex:04-three-components}
设 \(k\) 为域。证明 \(I=(x^2y,xy^2)\subset k[x,y]\) 有极小准素分解
\[
I=(x)\cap(y)\cap(x,y)^3.
\]
求关联素理想，并给出一个零化子为嵌入素理想的元素。
\end{exercise}
\begin{solution}
前两个理想的交为 \((xy)\)，其中全次数至少三的单项式恰被 \(x^2y\) 或 \(xy^2\) 整除，故再与 \((x,y)^3\) 取交得到 \(I\)。第三因子因根极大而准素。删去第三因子会增加 \(xy\)，删去第一因子会增加 \(y^3\)，删去第二因子会增加 \(x^3\)，因此分解极小。关联素理想为 \((x),(y),(x,y)\)，最后一个嵌入。非零元 \(\overline{xy}\) 的零化子为 \((x,y)\)：若 \(fxy\in(x^2y,xy^2)\)，可写成 \(fxy=xy(xa+yb)\)，在整环 \(k[x,y]\) 中约去非零元 \(xy\) 得 \(f=xa+yb\in(x,y)\)；反向直接相乘即可。
\end{solution}

\begin{exercise}\label{b3:ex:04-saturation-components}
设 \(k\) 为域。对 \(I=(x^2,xy)\subset R=k[x,y]\)，计算 \(I:y^\infty\) 与 \(I:x^\infty\)，并用准素分解解释被保留与被消去的准素分量。
\end{exercise}
\begin{solution}
对每个 \(n\ge1\)，\(fy^n\in I\) 当且仅当 \(f\in(x)\)，故 \(I:y^\infty=(x)\)。又 \(x^2\in I\)，所以 \(I:x^\infty=R\)。分解 \(I=(x)\cap(x^2,y)\) 中，倒置 \(y\) 保留根为 \((x)\) 的孤立准素分量，消去根含 \(y\) 的嵌入准素分量；倒置 \(x\) 则同时消去两个准素分量。饱和是相应局部化扩张的收缩，因此与直接计算一致。
\end{solution}

\begin{exercise}\label{b3:ex:04-maximal-components-crt}
设 \(R\) 为交换含幺 Noether 环，\(I\subsetneq R\) 为真理想，且 \(R/I\) 的全部关联素理想都是极大理想。证明 \(I\) 的极小准素分解除分量排列外唯一；若记作 \(I=\bigcap_{i=1}^r Q_i\)，则 \(R/I\cong\prod_{i=1}^rR/Q_i\)。
\end{exercise}
\begin{solution}
由\cref{b3:thm:primary-associated-primes}，极小准素分解的根素理想恰为 \(\operatorname{Ass}_R(R/I)\)。不同极大理想互不包含，所以这些关联素理想都孤立，\cref{b3:thm:primary-isolated-uniqueness}使各对应准素分量唯一。因此分解仅差分量的排列。若 \(Q_i+Q_j\ne R\)，取包含此和的极大理想 \(\mathfrak m\)，它包含 \(\sqrt{Q_i}\) 与 \(\sqrt{Q_j}\)，故两者都等于 \(\mathfrak m\)，与根互异矛盾。因此各 \(Q_i\) 两两互素，中国剩余定理给出所述乘积分解，其核正是 \(\bigcap_iQ_i=I\)。
\end{solution}

\begin{exercise}\label{b3:ex:04-parameter-embedded-components}
设 \(k\) 为域。对任意 \(c\in k\)，令 \(Q_c=(x^2,y-cx)\subset R=k[x,y]\)。证明 \((x)\cap Q_c=(x^2,xy)\)，并证明 \(c\ne d\) 时 \(Q_c\ne Q_d\)。
\end{exercise}
\begin{solution}
\(R/Q_c\cong k[x]/(x^2)\)，其中 \(y=cx\)，所以 \(Q_c\) 的根为 \((x,y)\)，它准素。若 \(xf\in Q_c\)，在此商环中 \(xf=x\cdot f(0,0)\)，而 \(x\ne0\)，故 \(f(0,0)=0\)，即 \(f\in(x,y)\)。因此交包含于 \((x^2,xy)\)，反向由 \(xy=x(y-cx)+cx^2\) 立即验证。若 \(c\ne d\)，\(y-cx\) 在 \(R/Q_d\) 中等于 \((d-c)x\ne0\)，所以不属于 \(Q_d\)。这给出另一类明确不同的嵌入准素分量。
\end{solution}

\begin{exercise}\label{b3:ex:04-diagonal-presentation}
设 \(k\) 为域，\(R=k[t]\)、\(M=R^2\)，\(N\) 是矩阵 \(\operatorname{diag}\bigl(t^2,(t-1)^3\bigr)\) 的像。给出 \(N\) 的极小准素分解、\(M/N\) 的关联素理想，以及商模的 \(k\) 维数。
\end{exercise}
\begin{solution}
令 \(N_0=t^2Re_1\oplus Re_2\)、\(N_1=Re_1\oplus(t-1)^3Re_2\)，则逐坐标取交得 \(N=N_0\cap N_1\)。两个商模分别为 \(R/(t^2)\) 与 \(R/\bigl((t-1)^3\bigr)\)，因而两因子分别属于 \((t)\)、\((t-1)\)，且无冗余。关联素理想就是这两个极大理想。商模同构于上述两个循环模的直和，维数为 \(2+3=5\)。
\end{solution}

\begin{exercise}\label{b3:ex:04-finiteness-primary-module}
设 \(k\) 为域，令 \(R=k[x]\)、\(P=\prod_{n\ge1}R/(x^n)\)。计算 \(\operatorname{Ann}_R(P)\)、\(\operatorname{Ass}_R(P)\) 与 \(Z_R(P)\)，并与\cref{b3:prop:primary-submodule-associated}证明中的无限直和反例比较。说明为什么每个坐标因子都被 \(x\) 的某个幂零化，直积里仍能出现零化子为 \((0)\) 的元素。
\end{exercise}
\begin{solution}
\(\operatorname{Ann}_R(P)=\bigcap_{n\ge1}(x^n)=(0)\)。对任意非零坐标元，它的零化子为某个 \((x^r)\)、\(r\ge1\)。一个非零序列的零化子是其非零坐标零化子的交：若这些指数有界，就取其中最大值，得到 \((x^r)\)；若指数无界，交为 \((0)\)。所以素零化子只能是 \((x)\) 或 \((0)\)。仅第一个坐标为 \(\bar1\) 的序列实现 \((x)\)，而所有坐标都为 \(\bar1\) 的序列实现 \((0)\)，故
\[
\operatorname{Ass}_R(P)=\{(0),(x)\}.
\]
若 \(f\notin(x)\)，它在每个 \(R/(x^n)\) 中都是单位，因此乘 \(f\) 在直积上也是单射；若 \(f\in(x)\)，它会零化仅第一个坐标非零的序列。所以 \(Z_R(P)=(x)\)，关联素理想 \((0)\) 没有向此并集增加元素。

无限直和中的元素只用有限个坐标，每个元素仍被某个 \(x^r\) 零化；无限直积没有这种限制，序列 \((\bar1,\bar1,\ldots)\) 的坐标需要无界的幂，故没有非零多项式能同时零化它。两个模本身的零化子都为 \((0)\)，零因子集合都为 \((x)\)，关联素理想集合却不同。
\end{solution}

\begin{exercise}\label{b3:ex:04-primary-not-irreducible-module}
设 \(R\) 为交换含幺环，\(\mathfrak p\) 为 \(R\) 的素理想，\(M=R/\mathfrak p\oplus R/\mathfrak p\)。直接证明零子模准素但可约。
\end{exercise}
\begin{solution}
若 \(a(m_1,m_2)=0\) 且元组非零，至少一个坐标在整环 \(R/\mathfrak p\) 中非零，故 \(a\in\mathfrak p\)。于是 \(aM=0\)，准素条件甚至以指数一成立。另一方面，两个坐标子模 \(R/\mathfrak p\oplus0\) 与 \(0\oplus R/\mathfrak p\) 都非零，交为零，所以零子模可约。子模准素与不能分解为交是不同性质。
\end{solution}

\begin{exercise}\label{b3:ex:04-triangular-presentation}
设 \(k\) 为域，在 \(R=k[x,y]\) 中令
\[
\Phi=\begin{pmatrix}x&y\\0&x\end{pmatrix},\qquad
N=\operatorname{im}\Phi\subset R^2.
\]
证明 \(R^2/N\) 的关联素理想只有 \((x)\)，因而 \(N\) 是准素子模。
\end{exercise}
\begin{solution}
记 \(\Psi=\begin{pmatrix}x&-y\\0&x\end{pmatrix}\)，则 \(\Phi\Psi=\Psi\Phi=x^2I_2\)。映射 \(v\mapsto\Psi v\bmod x^2R^2\) 诱导
\[
R^2/N\longrightarrow\bigl(R/(x^2)\bigr)^2.
\]
它是单射：若 \(\Psi v=x^2w\)，乘以 \(\Phi\) 得 \(x^2v=x^2\Phi w\)，而 \(R\) 是整环，所以 \(v=\Phi w\)。目标的关联素理想只有 \((x)\)，故源的关联集合包含于它。源非零，例如模掉 \((x,y)\) 后矩阵为零而余核为 \(k^2\)，所以关联素理想存在，必须恰为 \((x)\)。再用\cref{b3:prop:primary-submodule-associated}得 \(N\) 准素。
\end{solution}

\begin{exercise}\label{b3:ex:04-three-axes}
设 \(k\) 为域，在 \(R=k[x,y,z]\) 中令 \(I=(xy,xz,yz)\)。证明 \(I\) 为根理想，求 \(R/I\) 的关联素理想，并判断 \(\overline{x+y+z}\) 是否为零因子。
\end{exercise}
\begin{solution}
由单项式整除判据，\(I=(x,y)\cap(x,z)\cap(y,z)\)。它是素理想的交，故为根理想；三个因子互不包含且无冗余，因此关联素理想为这三个理想，或在商环中写作它们的像。\(x+y+z\) 模这三个理想分别为 \(z,y,x\)，均非零，故不在任何关联素理想中。由\cref{b3:thm:associated-zero-divisors}，其像为非零因子。
\end{solution}

\begin{exercise}\label{b3:ex:04-embedded-thickness}
设 \(k\) 为域，\(r\ge1\)，令 \(A=k[x,y]/(x^2,xy^r)\)。求幂零理想的 \(k\) 维数、它的支集，以及 \(A\) 的关联素理想。
\end{exercise}
\begin{solution}
根为 \((x)\)，所以幂零理想为 \((\bar x)\)，其 \(k\) 基为 \(\bar x,\bar x\bar y,\ldots,\bar x\bar y^{r-1}\)，维数 \(r\)。零化子为 \((\bar x,\bar y^r)\)，故支集只有 \((\bar x,\bar y)\) 这一点。直接验证
\[
(x^2,xy^r)=(x)\cap(x^2,y^r)
\]
给出极小准素分解，两个关联素理想为 \((\bar x)\) 与 \((\bar x,\bar y)\)。后者也由元素 \(\bar x\bar y^{r-1}\) 的零化子实现。倒置 \(y\) 后幂零理想消失，而不同的 \(r\) 给出不同维数的原点幂零部分。
\end{solution}

\begin{exercise}\label{b3:ex:04-annihilator-radical}
设 \(R\) 为交换含幺 Noether 环，\(M\ne0\) 为有限生成 \(R\) 模，\(\operatorname{Ass}_R(M)=\{\mathfrak p_1,\ldots,\mathfrak p_r\}\)。证明
\[
\sqrt{\operatorname{Ann}_R(M)}=\bigcap_{i=1}^r\mathfrak p_i,
\]
并证明某个 \(n\ge1\) 满足 \((\mathfrak p_1\cdots\mathfrak p_r)^nM=0\)。
\end{exercise}
\begin{solution}
关联素理想集合非空且有限，所以每个 \(\mathfrak p\in\operatorname{Ass}_R(M)\) 都包含某个极小关联素理想 \(\mathfrak q\)。因此，属于全部极小关联素理想的元素也属于每个关联素理想；反向包含则显然。由\cref{b3:prop:minimal-primes-associated}及\cref{b3:thm:radical-prime-intersection}，得到
\[
\begin{aligned}
\bigcap_{\mathfrak p\in\operatorname{Ass}_R(M)}\mathfrak p
&=\bigcap_{\mathfrak q\in\min\operatorname{Ass}_R(M)}\mathfrak q\\
&=\bigcap_{\mathfrak q\in\min V(\operatorname{Ann}_R(M))}\mathfrak q
=\sqrt{\operatorname{Ann}_R(M)}.
\end{aligned}
\]
记这个交为 \(J=(a_1,\ldots,a_s)\)，有限生成来自 Noether 性。对每个 \(a_i\) 选 \(n_i\ge1\) 使 \(a_i^{n_i}M=0\)。令 \(n=1+\sum_i(n_i-1)\)，则任意 \(n\) 次生成元乘积中某个 \(a_i\) 至少出现 \(n_i\) 次，所以 \(J^nM=0\)。因为 \(\mathfrak p_1\cdots\mathfrak p_r\subseteq J\)，所需结论成立。若 \(J=0\)，直接取 \(n=1\)。
\end{solution}
